\begin{abstract}
We present a method for recovering exact closed-form solutions to partial
differential equations from data.  The core technical contributions are twofold.
First, a \emph{normalization pipeline} that resolves gauge freedoms in the
parameterization of composed elementary functions---absorbing exponential biases
into multiplicative coefficients and canonicalizing trigonometric
phases---enabling reliable snapping of learned parameters to exact symbolic
values.  Second, a \emph{single-term search strategy} that trains each candidate
expression in isolation, avoiding the signal-splitting failure mode of
multi-term models with sparsity regularization.

These techniques operate within a structured ansatz we call the
\textsc{ComposeHead}, which represents candidate solutions as products of
axis-aligned elementary functions---matching the separable structure common to
analytical PDE solutions.  The primitive basis
$\{\sin, \cos, \exp, \ln\}$ is grounded in the EML (Exp-Minus-Log) universality
theorem of Odrzywo{\l}ek~\cite{odrzywołek2026}, which provides a principled
justification for this function library and guarantees closure under
differentiation---a property essential for computing PDE residuals within the
compositional graph.

We demonstrate the approach on the two-dimensional Navier--Stokes equations,
recovering the Taylor--Green vortex solution
$u = \sin(x)\cos(y)\exp(-t/5)$,\;$v = -\cos(x)\sin(y)\exp(-t/5)$
at machine precision (pointwise error ${<}\,10^{-6}$).  Ablation studies confirm
that normalization is the most impactful design choice, improving snap accuracy
by five orders of magnitude.
\end{abstract}
